\section{Agent harness details}
\label{app:harness}
\paragraph{Sandbox.} Agent-written code and the replay of every final program run in a sandbox with no network and writes confined to the episode folder: bubblewrap where available; otherwise an unprivileged user namespace in which every mount is remounted read-only and verified (the sandbox refuses to run if any mount stays writable), temporary directories are replaced by empty file systems, the episode folder is bound back read-write, and all capabilities are dropped with no-new-privileges set.

\paragraph{Context.} The last six tool results are kept in full; older results, older assistant text and older code arguments are truncated to 300 characters in what is sent back. Tool results are capped at 8{,}000 characters while remaining valid JSON. On a context-length error the kept window halves.

\paragraph{Models.}
\begin{center}
\scriptsize
\resizebox{\linewidth}{!}{% generated by bench/paper_tables.py -- do not edit
\begin{tabular}{@{}llp{6.6cm}@{}}
\toprule
model & temperature & request extras and vLLM parsers \\
\midrule
Gemma-4-31B (\texttt{google/gemma-4-31B-it}) & 1.0 & \texttt{\{"top\_p": 0.95, "top\_k": 64, "chat\_template\_kwargs": \{"enable\_thinking": true\}\}}; \texttt{--reasoning-parser gemma4 --tool-call-parser gemma4} \\
Qwen3.6-35B-A3B (\texttt{Qwen/Qwen3.6-35B-A3B-FP8}) & 1.0 & \texttt{\{"top\_p": 0.95, "top\_k": 20, "min\_p": 0.0, "presence\_penalty": 0.0, "chat\_template\_kwargs": \{"enable\_thinking": true\}\}}; \texttt{--reasoning-parser qwen3 --tool-call-parser qwen3\_coder --language-model-only} \\
gpt-oss-120b (\texttt{openai/gpt-oss-120b}) & 1.0 & \texttt{\{"top\_p": 1.0\}}; \texttt{--reasoning-parser openai\_gptoss --tool-call-parser openai} \\
Devstral-S-2-24B (\texttt{mistralai/Devstral-Small-2-24B-Instruct-2512}) & 0.15 & \texttt{\{\}}; \texttt{--tool-call-parser mistral --tokenizer-mode mistral --config-format mistral --load-format mistral} \\
gpt-oss-20b (\texttt{openai/gpt-oss-20b}) & 1.0 & \texttt{\{"top\_p": 1.0\}}; \texttt{--reasoning-parser openai\_gptoss --tool-call-parser openai} \\
Qwen3.8-27B (\texttt{Qwen/Qwen3.8-27B-FP8}) & 1.0 & \texttt{\{"top\_p": 0.95, "top\_k": 20, "min\_p": 0.0, "presence\_penalty": 0.0, "chat\_template\_kwargs": \{"enable\_thinking": true\}\}}; \texttt{--reasoning-parser qwen3 --tool-call-parser qwen3\_xml --language-model-only} \\
\bottomrule
\end{tabular}
}
\end{center}

\paragraph{Prompts.} The three system prompts differ only in the tool paragraph and the strategy phrase.
% generated by bench/paper_tables.py -- do not edit
\paragraph{System prompt, \texttt{code} toolset.}
{\scriptsize\begin{verbatim}
You are a CAD engineer. Your job: build a solid CAD model that matches a target 3D part as
    closely as possible, measured by volumetric IoU (intersection over union) against the
    target.

You cannot see the target. You learn about it only through tools:
- describe_target: size, volume, and the flat faces (directions and plane positions).
- section: exact cross-section polygons of the target (and of your model) in any
    axis-aligned plane.
- add_code: append a solid built by your CadQuery code to the model as a union ('add') or a
    difference ('cut').
- evaluate: current IoU and where material is missing or extra.
- remove_step, list_steps, finish.

Coordinates: everything is in the normalized frame (the target is centered and scaled so
    its largest extent spans [-1, 1]). Your CadQuery code (cq and math are imported) must
    assign the final solid to `result`, e.g.
    result = cq.Workplane('XY').box(1.2, 0.8, 0.1).translate((0, 0, 0.35))
    result = cq.Workplane('XY', origin=(0, 0, -0.2)).circle(0.3).extrude(0.4)
    result =
    cq.Workplane('XZ').polyline([(0,0),(0.5,0),(0.5,0.2),(0,0.2)]).close().revolve(360,
    (0,0,0), (0,1,0))

Good strategy: describe the target, read sections to get exact profile shapes, build the
    main body first, then add features and cut holes/pockets, checking evaluate after
    changes and removing steps that lower IoU. The best model seen at any point is kept, so
    experimenting is safe. Work efficiently: you have a limited number of tool calls. Call
    finish when you cannot improve further.
\end{verbatim}}
\paragraph{System prompt, \texttt{codesec} toolset.}
{\scriptsize\begin{verbatim}
You are a CAD engineer. Your job: build a solid CAD model that matches a target 3D part as
    closely as possible, measured by volumetric IoU (intersection over union) against the
    target.

You cannot see the target. You learn about it only through tools:
- describe_target: size, volume, and the flat faces (directions and plane positions).
- section: exact cross-section polygons of the target (and of your model) in any
    axis-aligned plane.
- section_extrude: a CADFit stage that extrudes the target's whole cross-section at a plane
    through the full extent of the part. Each solid it makes becomes a scored candidate:
    list_candidates / show_candidate show them; precision near 1 means the candidate lies
    inside the target.
- add_candidate / add_code: append a solid to the model as a union ('add') or a difference
    ('cut').
- evaluate: current IoU and where material is missing or extra.
- remove_step, list_steps, finish.

Coordinates: everything is in the normalized frame (the target is centered and scaled so
    its largest extent spans [-1, 1]). Your CadQuery code (cq and math are imported) must
    assign the final solid to `result`, e.g.
    result = cq.Workplane('XY').box(1.2, 0.8, 0.1).translate((0, 0, 0.35))
    result = cq.Workplane('XY', origin=(0, 0, -0.2)).circle(0.3).extrude(0.4)
    result =
    cq.Workplane('XZ').polyline([(0,0),(0.5,0),(0.5,0.2),(0,0.2)]).close().revolve(360,
    (0,0,0), (0,1,0))

Good strategy: describe the target, read sections to get exact profile shapes and extrude
    whole cross-sections, build the main body first, then add features and cut
    holes/pockets, checking evaluate after changes and removing steps that lower IoU. The
    best model seen at any point is kept, so experimenting is safe. Work efficiently: you
    have a limited number of tool calls. Call finish when you cannot improve further.
\end{verbatim}}
\paragraph{System prompt, \texttt{skills} toolset.}
{\scriptsize\begin{verbatim}
You are a CAD engineer. Your job: build a solid CAD model that matches a target 3D part as
    closely as possible, measured by volumetric IoU (intersection over union) against the
    target.

You cannot see the target. You learn about it only through tools:
- describe_target: size, volume, and the flat faces (directions and plane positions).
- section: exact cross-section polygons of the target (and of your model) in any
    axis-aligned plane.
- CADFit skills, the stages of a CAD reconstruction pipeline, which you run yourself:
    detect_directions (likely extrusion directions), extract_sketch (slice the target at a
    plane into sketch profiles fitted with lines and arcs), fit_extrude / fit_revolve (turn
    a sketch into extrusion / revolve solids that match the target), section_extrude
    (extrude a whole cross-section through the part). Each solid they make becomes a scored
    candidate: list_candidates / show_candidate show them; precision near 1 means the
    candidate lies inside the target.
- add_candidate / add_code: append a solid to the model as a union ('add') or a difference
    ('cut').
- evaluate: current IoU and where material is missing or extra.
- remove_step, list_steps, finish.

Coordinates: everything is in the normalized frame (the target is centered and scaled so
    its largest extent spans [-1, 1]). Your CadQuery code (cq and math are imported) must
    assign the final solid to `result`, e.g.
    result = cq.Workplane('XY').box(1.2, 0.8, 0.1).translate((0, 0, 0.35))
    result = cq.Workplane('XY', origin=(0, 0, -0.2)).circle(0.3).extrude(0.4)
    result =
    cq.Workplane('XZ').polyline([(0,0),(0.5,0),(0.5,0.2),(0,0.2)]).close().revolve(360,
    (0,0,0), (0,1,0))

Good strategy: describe the target, extract sketches on its main planes and fit them into
    candidates, build the main body first, then add features and cut holes/pockets,
    checking evaluate after changes and removing steps that lower IoU. The best model seen
    at any point is kept, so experimenting is safe. Work efficiently: you have a limited
    number of tool calls. Call finish when you cannot improve further.
\end{verbatim}}


\section{Tool use}
\label{app:tools}
Tool calls per episode in the skills toolset, over the \nAgentParts{} episodes of each model (\texttt{results/tool\_usage.json}).
\begin{center}
\scriptsize
\resizebox{\linewidth}{!}{% generated by bench/paper_tables.py -- do not edit
\begin{tabular}{@{}lcccccccccccc@{}}
\toprule
model & calls & \rotatebox{70}{\texttt{extract\_sketch}} & \rotatebox{70}{\texttt{fit\_extrude}} & \rotatebox{70}{\texttt{fit\_revolve}} & \rotatebox{70}{\texttt{section\_extrude}} & \rotatebox{70}{\texttt{detect\_directions}} & \rotatebox{70}{\texttt{section}} & \rotatebox{70}{\texttt{add\_candidate}} & \rotatebox{70}{\texttt{add\_code}} & \rotatebox{70}{\texttt{evaluate}} & \rotatebox{70}{\texttt{remove\_step}} & errors \\
\midrule
Gemma-4-31B & 27.9 & 3.1 & 2.7 & 0.1 & 1.1 & 0.4 & 6.0 & 2.9 & 3.5 & 1.9 & 2.6 & 17/559 \\
Qwen3.6-35B-A3B & 37.2 & 4.9 & 2.2 & 0.0 & 1.6 & 1.0 & 12.3 & 2.5 & 3.4 & 2.2 & 2.5 & 47/744 \\
gpt-oss-120b & 38.2 & 4.8 & 3.6 & 0.1 & 1.6 & 1.0 & 6.6 & 4.0 & 3.6 & 3.0 & 3.4 & 22/764 \\
Devstral-S-2-24B & 37.2 & 3.8 & 6.8 & 0.3 & 1.2 & 0.9 & 2.7 & 7.2 & 3.4 & 4.1 & 3.9 & 12/745 \\
gpt-oss-20b & 25.1 & 2.1 & 1.8 & 0.1 & 1.6 & 0.9 & 6.2 & 2.6 & 2.2 & 2.6 & 1.2 & 44/503 \\
Qwen3.8-27B & 22.2 & 1.8 & 1.6 & 0.0 & 0.3 & 1.0 & 7.9 & 1.3 & 3.1 & 1.3 & 1.8 & 27/445 \\
\bottomrule
\end{tabular}
}
\end{center}

\section{Qwen3.8 token-limit failures and a post-hoc sensitivity run}
\label{app:qwen}
The harness sends ``Your reply hit the token limit before you called a tool. Think briefly, then call a tool now.'' only when a reply ends because it reached the 16{,}384-token limit; after two consecutive nudges without a tool call the episode stops. In all \qwenTraceN{} Qwen3.8 code episodes that stopped this way, the last nudge was the token-limit message, and the final reply was \qwenTraceMinK--\qwenTraceMaxK{} thousand characters of reasoning with no content (\texttt{results/qwen38\_traces.json}). After seeing these outcomes we reran Qwen3.8 on the same parts with 32{,}768 tokens per turn. It scored \qwenPosthocSkills{} with skills and \qwenPosthocCode{} with code (versus \mQwenEightSkills{} and \mQwenEightCode{} at 16{,}384). Because this run was decided after seeing results, it is excluded from \cref{tab:agents}; most of its episodes ended at the 30-minute limit instead.

\section{Refining pipeline output with an LLM agent}
\label{app:refine}
In an earlier study, Qwen3.8-27B (INT4) refined arm C's output on the \refN{} fresh parts whose C mesh IoU was below 0.90 (gate fixed before running), with access to C's candidate pool, 40 calls and 20 minutes, in two runs. Mean mesh IoU: pipeline \refPipe; agent \refAgentOne{} and \refAgentTwo{} (wins--losses against the pipeline at a $\pm 0.01$ threshold: \refAgentOneWL{} and \refAgentTwoWL); a no-op agent that only re-exports the pipeline model \refNoop; an equal-time scripted search over the same pool \refSearch{} (\refSearchWL). The mean gain comes from a few parts and does not repeat consistently. Per-part records were not retained, so no intervals or executed-program scores are available for this study.

\section{Per-part results}
\label{app:perpart}
{\scriptsize
% generated by bench/paper_tables.py -- do not edit
\begin{longtable}{@{}llcccccccc@{}}
\toprule
part & split & A mesh & A CAD & B mesh & B CAD & C mesh & C CAD & D mesh & D CAD \\
\midrule\endhead
00038011 & A & 0.844 & 0.844 & 0.939 & 0.937 & 0.940 & 0.940 & 0.940 & 0.940 \\
00046097 & A & 0.896 & 0.915 & 0.939 & 0.940 & 0.942 & 0.943 & 0.942 & 0.943 \\
00053172 & E & 0.994 & 0.994 & 0.994 & 0.994 & 0.994 & 0.994 & 0.998 & 0.998 \\
00090768 & A & 0.962 & 0.962 & 0.971 & 0.971 & 0.999 & 0.967 & 0.999 & 0.002 \\
00095369 & A & 0.882 & 0.882 & 0.893 & 0.893 & 0.893 & 0.893 & 0.948 & 0.948 \\
00101419 & A & 0.965 & 0.965 & 0.978 & 0.978 & 0.978 & 0.978 & 1.000 & 1.000 \\
00140273 & E & 0.884 & 0.884 & 0.884 & 0.884 & 0.884 & 0.884 & 0.894 & 0.894 \\
00153120 & E & 0.771 & 0.771 & 0.857 & 0.857 & 0.857 & 0.857 & 0.996 & 0.996 \\
00154184 & E & 0.000 & --$^\dagger$ & 0.980 & 0.981 & 0.984 & 0.984 & 0.984 & 0.984 \\
00161800 & A & 0.985 & 0.985 & 0.985 & 0.985 & 0.985 & 0.985 & 0.985 & 0.985 \\
00181393 & A & 0.863 & 0.863 & 0.863 & 0.863 & 0.863 & 0.863 & 0.863 & 0.863 \\
00208290 & E & 0.989 & 0.989 & 0.989 & 0.989 & 0.989 & 0.989 & 0.992 & 0.992 \\
00214289 & A & 0.997 & 0.997 & 0.989 & 0.989 & 0.989 & 0.989 & 0.989 & 0.989 \\
00232213 & E & 0.837 & 0.837 & 0.869 & 0.000$^\dagger$ & 0.861 & 0.873$^*$ & 0.869 & 0.000$^\dagger$ \\
00246395 & E & 0.959 & 0.959 & 0.959 & 0.959 & 0.959 & 0.959 & 0.997 & 0.997 \\
00247424 & E & 0.000 & --$^\dagger$ & 0.000 & --$^\dagger$ & 0.000 & --$^\dagger$ & 0.000 & --$^\dagger$ \\
00249450 & E & 0.926 & 0.926 & 0.960 & 0.960 & 0.960 & 0.960 & 0.995 & 0.995 \\
00256926 & E & 0.814 & 0.814 & 0.892 & 0.892 & 0.899 & 0.902 & 1.000 & 1.000 \\
00266468 & E & 0.723 & 0.724 & 0.746 & 0.746 & 0.749 & 0.746 & 1.000 & 1.000 \\
00273535 & E & 0.912 & 0.912 & 0.995 & 0.995 & 0.995 & 0.995 & 1.000 & 1.000 \\
00299355 & E & 0.975 & 0.976 & 0.975 & 0.976 & 0.975 & 0.976 & 0.990 & 0.990 \\
00305871 & E & 0.960 & 0.960 & 0.961 & 0.961 & 0.961 & 0.961 & 1.000 & 1.000 \\
00312587 & E & 0.973 & 0.973 & 0.975 & 0.974 & 0.975 & 0.974 & 1.000 & 1.000 \\
00317970 & E & 0.182 & 0.182 & 0.733 & 0.743 & 0.746 & 0.747 & 0.993 & 0.993 \\
00354176 & A & 0.896 & 0.896 & 0.903 & 0.904 & 0.903 & 0.904 & 0.903 & 0.904 \\
00365747 & A & 0.848 & 0.848 & 0.880 & 0.000$^\dagger$ & 0.888 & 0.000$^\dagger$ & 0.961 & 0.956$^*$ \\
00382121 & A & 0.920 & 0.920 & 0.933 & 0.933 & 0.936 & 0.934 & 0.936 & 0.936 \\
00492401 & A & 0.975 & 0.986 & 0.986 & 0.986 & 0.987 & 0.988 & 0.987 & 0.987 \\
00528181 & A & 0.860 & 0.876 & 0.915 & 0.916$^*$ & 0.882 & 0.882 & 0.960 & 0.958 \\
00562677 & A & 0.867 & 0.867 & 0.889 & 0.889 & 0.889 & 0.889 & 0.889 & 0.889 \\
00589302 & A & 0.921 & 0.921 & 0.966 & 0.966 & 0.969 & 0.937$^*$ & 0.949 & 0.949 \\
00630878 & A & 0.848 & 0.849 & 0.942 & 0.944 & 0.944 & 0.944 & 0.944 & 0.944 \\
00634093 & A & 0.931 & 0.931 & 0.941 & 0.942 & 0.946 & 0.946 & 0.946 & 0.946 \\
00642725 & A & 0.908 & 0.908 & 0.931 & 0.000$^\dagger$ & 0.930 & 0.000$^\dagger$ & 0.936 & 0.000$^\dagger$ \\
00657371 & E & 0.901 & 0.902 & 0.901 & 0.902 & 0.902 & 0.903 & 0.987 & 0.987 \\
00678171 & E & 0.935 & 0.935 & 0.945 & 0.945 & 0.945 & 0.945 & 0.949 & 0.949 \\
00700032 & A & 0.855 & 0.855 & 0.885 & 0.881 & 0.886 & 0.886 & 0.955 & 0.955 \\
00709520 & A & 0.888 & 0.889 & 0.888 & 0.889 & 0.888 & 0.889 & 0.888 & 0.889 \\
00729452 & A & 0.051 & 0.053 & 0.280 & 0.281 & 0.280 & 0.281 & 0.280 & 0.281 \\
00730725 & E & 0.018 & 0.018 & 0.018 & 0.018 & 0.018 & 0.018 & 0.949 & 0.949 \\
00771160 & E & 0.890 & 0.900 & 0.974 & 0.974 & 0.985 & 0.982$^*$ & 0.985 & 0.982$^*$ \\
00771600 & E & 0.919 & 0.919 & 0.919 & 0.919 & 0.920 & 0.920 & 0.999 & 0.999 \\
00780711 & E & 0.911 & 0.911 & 0.973 & 0.973 & 0.983 & 0.983 & 0.983 & 0.983 \\
00791454 & E & 0.970 & 0.970 & 0.970 & 0.970 & 0.970 & 0.970 & 0.990 & 0.990 \\
00795057 & E & 0.949 & 0.949 & 0.969 & 0.969 & 0.980 & 0.980 & 0.997 & 0.997 \\
00814797 & E & 0.933 & 0.933 & 0.970 & 0.970 & 0.970 & 0.970 & 0.999 & 0.999 \\
00834538 & A & 0.920 & 0.920 & 0.948 & 0.948 & 0.954 & 0.949 & 0.968 & 0.968 \\
00846352 & A & 0.859 & 0.859 & 0.859 & 0.859 & 0.867 & 0.867 & 0.867 & 0.867 \\
00862084 & A & 0.593 & 0.593 & 0.553 & 0.557 & 0.639 & 0.648 & 0.999 & 0.999 \\
00862085 & E & 0.422 & 0.422 & 0.475 & 0.475 & 0.475 & 0.475 & 0.999 & 0.999 \\
00862953 & E & 0.796 & 0.796 & 0.802 & 0.804 & 0.802 & 0.804 & 0.980 & 0.980 \\
00883892 & E & 0.792 & 0.792 & 0.905 & 0.906 & 0.913 & 0.000$^\dagger$ & 0.900 & 0.900 \\
00892314 & A & 0.916 & 0.916 & 0.933 & 0.940$^*$ & 0.935 & 0.000$^\dagger$ & 0.959 & 0.959 \\
00904252 & A & 0.977 & 0.977 & 0.977 & 0.977 & 0.977 & 0.977 & 0.977 & 0.977 \\
00941601 & A & 0.923 & 0.923 & 0.943 & 0.943 & 0.947 & 0.947 & 0.949 & 0.949 \\
00953658 & E & 0.973 & 0.973 & 0.976 & 0.976 & 0.978 & 0.978 & 0.982 & 0.982 \\
00966049 & A & 0.815 & 0.815 & 0.945 & 0.818 & 0.945 & 0.818 & 0.950 & 0.821 \\
00966713 & E & 0.994 & 0.994 & 0.994 & 0.994 & 0.994 & 0.994 & 0.998 & 0.998 \\
00983172 & A & 0.877 & 0.877 & 0.993 & 0.985 & 0.993 & 0.000$^\dagger$ & 0.984 & 0.984 \\
00992810 & A & 0.879 & 0.879 & 0.995 & 0.995 & 0.995 & 0.995 & 0.997 & 0.997 \\
\bottomrule
\end{longtable}

}
$^\dagger$ program does not build; $^*$ builds but fails OpenCascade validity.

\begin{table}[h]
\centering
\caption{Per-part \iouc{} of the agent study (failures 0).}
\resizebox{\linewidth}{!}{% generated by bench/paper_tables.py -- do not edit
\begin{tabular}{@{}lcccccccccccccc@{}}
\toprule
part & \multicolumn{2}{c}{Gemma-4-31B} & \multicolumn{2}{c}{Qwen3.6-35B-A3B} & \multicolumn{2}{c}{gpt-oss-120b} & \multicolumn{2}{c}{Devstral-S-2-24B} & \multicolumn{2}{c}{gpt-oss-20b} & \multicolumn{2}{c}{Qwen3.8-27B} & policy & pipe D \\
 & code & skills & code & skills & code & skills & code & skills & code & skills & code & skills & & \\
\midrule
00046097 & 0.389 & 0.832 & 0.000 & 0.739 & 0.702 & 0.770 & 0.274 & 0.765 & 0.844 & 0.739 & 0.000 & 0.785 & 0.776 & 0.943 \\
00095369 & 0.955 & 0.951 & 0.945 & 0.872 & 0.866 & 0.953 & 0.086 & 0.838 & 0.873 & 0.872 & 0.945 & 0.986 & 0.899 & 0.948 \\
00140273 & 0.251 & 0.894 & 0.439 & 0.884 & 0.446 & 0.894 & 0.308 & 0.894 & 0.245 & 0.647 & 0.000 & 0.000 & 0.894 & 0.894 \\
00153120 & 0.976 & 0.996 & 0.410 & 0.996 & 0.976 & 0.996 & 0.227 & 0.971 & 0.579 & 0.996 & 0.749 & 0.996 & 0.996 & 0.996 \\
00154184 & 0.230 & 0.967 & 0.168 & 0.936 & 0.268 & 0.218 & 0.211 & 0.218 & 0.151 & 0.195 & 0.000 & 0.971 & 0.810 & 0.984 \\
00232213 & 0.121 & 0.279 & 0.399 & 0.304 & 0.228 & 0.276 & 0.126 & 0.458 & 0.243 & 0.000 & 0.000 & 0.000 & 0.403 & 0.000 \\
00249450 & 0.980 & 0.979 & 0.980 & 0.995 & 0.942 & 0.995 & 0.327 & 0.881 & 0.651 & 0.995 & 0.978 & 0.980 & 0.995 & 0.995 \\
00354176 & 0.021 & 0.929 & 0.221 & 0.929 & 0.313 & 0.930 & 0.221 & 0.929 & 0.281 & 0.929 & 0.566 & 0.929 & 0.000 & 0.904 \\
00492401 & 0.647 & 0.932 & 0.231 & 0.934 & 0.602 & 0.847 & 0.343 & 0.883 & 0.292 & 0.848 & 0.000 & 0.910 & 0.871 & 0.987 \\
00528181 & 0.221 & 0.889 & 0.472 & 0.765 & 0.196 & 0.889 & 0.252 & 0.706 & 0.164 & 0.765 & 0.000 & 0.732 & 0.784 & 0.958 \\
00589302 & 0.747 & 0.948 & 0.950 & 0.974 & 0.950 & 0.948 & 0.652 & 0.669 & 0.950 & 0.948 & 0.000 & 0.824 & 0.948 & 0.949 \\
00630878 & 0.252 & 0.951 & 0.000 & 0.658 & 0.255 & 0.747 & 0.181 & 0.676 & 0.265 & 0.615 & 0.000 & 0.835 & 0.746 & 0.944 \\
00700032 & 0.961 & 0.962 & 0.608 & 0.962 & 0.960 & 0.962 & 0.326 & 0.962 & 0.614 & 0.962 & 0.975 & 0.963 & 0.962 & 0.955 \\
00771600 & 0.944 & 0.999 & 0.618 & 0.999 & 0.618 & 0.621 & 0.713 & 0.919 & 0.000 & 0.999 & 0.000 & 0.985 & 0.999 & 0.999 \\
00834538 & 0.816 & 0.969 & 0.319 & 0.968 & 0.816 & 0.968 & 0.183 & 0.966 & 0.816 & 0.968 & 0.813 & 0.968 & 0.968 & 0.968 \\
00862084 & 0.933 & 0.999 & 0.157 & 0.999 & 0.489 & 0.905 & 0.161 & 0.798 & 0.124 & 0.993 & 0.926 & 0.975 & 0.999 & 0.999 \\
00883892 & 0.668 & 0.773 & 0.590 & 0.669 & 0.676 & 0.789 & 0.590 & 0.572 & 0.590 & 0.793 & 0.169 & 0.000 & 0.793 & 0.900 \\
00941601 & 0.289 & 0.898 & 0.332 & 0.740 & 0.030 & 0.873 & 0.489 & 0.731 & 0.338 & 0.000 & 0.000 & 0.916 & 0.740 & 0.949 \\
00953658 & 0.628 & 0.929 & 0.622 & 0.929 & 0.623 & 0.929 & 0.415 & 0.893 & 0.415 & 0.622 & 0.000 & 0.844 & 0.929 & 0.982 \\
00966049 & 0.830 & 0.764 & 0.836 & 0.823 & 0.176 & 0.735 & 0.011 & 0.790 & 0.000 & 0.763 & 0.018 & 0.000 & 0.679 & 0.821 \\
\bottomrule
\end{tabular}
}
\end{table}
